% Zumdahl Chapter 2 -- Atoms, Molecules, and Ions. Chapter-track guided notes.
% Follows the BOOK's sequence (not the CED's) -- see Chapters/README.md.
\documentclass{apchem-notes}

\apcunitword{Chapter}
\apcunit{2}
\apcunittitle{Atoms, Molecules, and Ions}
\apcdoctitle{Guided Notes}
\apcsubtitle{Zumdahl \S2.1--2.8 \textbullet\ PDF pp.~69--109 \textbullet\ 3 blocks}

\begin{document}
\apctitleblock

\begin{apcnote}[How this chapter fits the AP course]
Most of this chapter is \emph{foundation}, not tested CED content: the
historical laws and early experiments build the reasoning AP assesses in
Unit 1, and \S2.8 nomenclature is a prerequisite skill assumed silently by
every later unit. Sections 2.6--2.8 carry the durable payload --- ions,
formulas, and naming. Learn the names cold now and Units 4, 7, and 8 get
much easier.
\end{apcnote}

% =====================================================================
\blocklesson{From Laws to Atoms}{Zumdahl \S2.1--2.5}

\begin{objectives}
  \item State the three fundamental mass laws and explain how each supports
        atomic theory.
  \item Describe the experiments that revealed the electron and the nucleus,
        and what each one ruled out.
\end{objectives}

\reading{Zumdahl \S2.1--2.5}{69--82}

\begin{warmup}[3]
  \item How many protons in \ce{^{35}Cl}? \answerblank[2cm]{17}
  \item Which subatomic particle has negligible mass?
        \answerblank[2.5cm]{electron}
  \item Charge of a neutron: \answerblank[2cm]{zero}
\end{warmup}

\segment{INSTRUCTION A}{25}{The three mass laws}

\notesectionz{Laws that demanded atoms}{2.2}
\practice{6}

\renewcommand{\arraystretch}{1.5}
\noindent\begin{tabular}{@{}p{3.6cm}p{4.2cm}p{5.6cm}@{}}
\toprule
\textbf{Law} & \textbf{Statement} & \textbf{What it implies}\\
\midrule
Conservation of mass\newline(Lavoisier) &
  mass is neither \answerblank[1.8cm]{created} nor destroyed in a reaction &
  atoms are rearranged, not made or destroyed\\
Definite proportion\newline(Proust) &
  a given compound always has the \answerblank[2.4cm]{same} mass ratio of
  elements &
  compounds have fixed formulas\\
Multiple proportions\newline(Dalton) &
  masses of B combining with a fixed mass of A are in small
  \answerblank[3.0cm]{whole-number} ratios &
  atoms combine in discrete units --- the strongest argument\\
\bottomrule
\end{tabular}

\begin{workedexample}[Worked example: multiple proportions]
Two oxides of carbon. In compound I, 1.00 g C combines with 1.33 g O; in
compound II, 1.00 g C combines with 2.66 g O.
\[
  \frac{2.66}{1.33} = 2.00
\]
The oxygen masses are in a 2:1 ratio --- exactly what you'd expect if
compound I is \ce{CO} and II is \ce{CO2}. Fractional ratios would be
impossible if matter were continuous.
\end{workedexample}

\notesub{Dalton's atomic theory (1808)}

\begin{enumerate}[leftmargin=*,itemsep=0.15em]
  \item Elements are made of tiny particles called \answerblank[2.2cm]{atoms}.
  \item All atoms of a given element are \answerblank[2.4cm]{identical}
        --- \emph{later corrected by} \answerblank[2.2cm]{isotopes}.
  \item Atoms of different elements differ in
        \answerblank[3.4cm]{mass/properties}.
  \item Compounds form when atoms combine in fixed
        \answerblank[4.3cm]{whole-number ratios}.
  \item Chemical reactions only \answerblank[3cm]{rearrange} atoms.
\end{enumerate}

\segment{GUIDED PRACTICE}{15}{Law identification}

Which law does each observation illustrate?
\begin{enumerate}[leftmargin=*,itemsep=0.35em]
  \item Water from any source is 11.2\% H by mass.
        \answerblank[5.5cm]{definite proportion}
  \item \SI{10.0}{\gram} reactants yield \SI{10.0}{\gram} products.
        \answerblank[5.5cm]{conservation of mass}
  \item \ce{NO} and \ce{NO2}: O per gram of N is in a 1:2 ratio.
        \answerblank[5.5cm]{multiple proportions}
\end{enumerate}

\segment{INSTRUCTION B}{20}{Taking the atom apart}

\notesectionz{Three experiments, three conclusions}{2.4--2.5}
\practice{4}

\renewcommand{\arraystretch}{1.4}
\noindent\begin{tabular}{@{}p{3.4cm}p{4.4cm}p{5.6cm}@{}}
\toprule
\textbf{Experiment} & \textbf{Observation} & \textbf{Conclusion}\\
\midrule
Thomson\newline cathode ray &
  beam deflects toward the \answerblank[1.6cm]{positive} plate; same
  charge/mass for any gas &
  \answerblank[3.4cm]{electrons exist} in all atoms; ``plum pudding''\\
Millikan\newline oil drop &
  droplet charges are multiples of one value &
  charge of the \answerblank[2.2cm]{electron}; with $e/m$, gives its mass\\
Rutherford\newline gold foil &
  most $\alpha$ pass through; a few deflect \answerblank[2.7cm]{sharply back} &
  atom is mostly \answerblank[2.7cm]{empty space} with a tiny, dense,
  positive \answerblank[1.8cm]{nucleus}\\
\bottomrule
\end{tabular}

\begin{apctrap}
The gold-foil result is famous for what it \emph{disproved}. Plum pudding
predicted \emph{all} particles pass nearly straight through --- spread-out
positive charge can't repel hard. The rare large-angle backscatter is the
evidence, not the common straight-through path. AP asks for the anomaly.
\end{apctrap}

\segment{APPLICATION}{20}{Reasoning from evidence}

\begin{enumerate}[leftmargin=*,itemsep=0.4em]
  \item Thomson found the same charge-to-mass ratio no matter which gas
        filled the tube. Why did that matter?
        \answerlines[2]{The same particle came from every element ---
        electrons are a universal component of all atoms, not a property of
        one gas.}
  \item If the nucleus were the size of a marble, the atom would be about
        the size of a stadium. Which experimental result does that
        illustrate? \answerlines[2]{Rutherford's: nearly all $\alpha$
        particles passed undeflected because the atom is mostly empty
        space.}
\end{enumerate}

\begin{exitticket}
Dalton said all atoms of an element are identical. What later discovery
corrected this, and does it break the theory?
\answerlines[2]{Isotopes --- same protons, different neutrons. The theory
survives: atoms of an element still behave identically chemically, since
chemistry depends on electrons/protons, not neutrons.}
\end{exitticket}

% =====================================================================
\blocklesson{Atoms, Ions, and the Periodic Table}{Zumdahl \S2.5--2.7}

\begin{objectives}
  \item Use isotope notation to count protons, neutrons, and electrons.
  \item Predict common ion charges from periodic-table position.
\end{objectives}

\reading{Zumdahl \S2.5--2.7}{80--87}

\begin{warmup}[3]
  \item Which experiment found the nucleus? \answerblank[3cm]{gold foil}
  \item \ce{CO} vs \ce{CO2} illustrates which law?
        \answerblank[4cm]{multiple proportions}
  \item Sign of an electron's charge: \answerblank[2cm]{negative}
\end{warmup}

\segment{INSTRUCTION A}{25}{Isotope notation}

\notesectionz{Reading \ce{^{A}_{Z}X}}{2.5}
\practice{1}

\begin{itemize}[leftmargin=*,itemsep=0.2em]
  \item $Z$ = \term{atomic number} = number of
        \answerblank[2.2cm]{protons} --- fixes the element's identity.
  \item $A$ = \term{mass number} = protons $+$
        \answerblank[2.2cm]{neutrons}.
  \item Neutrons $= $ \answerblank[2cm]{$A - Z$}.
  \item Neutral atom: electrons $=$ \answerblank[2cm]{protons}.
        For an ion, adjust by the \answerblank[2cm]{charge}.
\end{itemize}

\begin{workedexample}[Worked example: three species]
\ce{^{56}_{26}Fe^{3+}}: 26 protons, $56-26 = 30$ neutrons,
$26-3 = 23$ electrons.\\
\ce{^{32}_{16}S^{2-}}: 16 protons, 16 neutrons, $16+2 = 18$ electrons.\\
\ce{^{40}_{18}Ar}: 18 protons, 22 neutrons, 18 electrons.\\
Note \ce{S^{2-}} and \ce{Ar} both have 18 electrons --- isoelectronic.
\end{workedexample}

\segment{GUIDED PRACTICE}{15}{Fill the table}

\renewcommand{\arraystretch}{1.5}
\noindent\begin{tabular}{@{}lcccc@{}}
\toprule
\textbf{Species} & \textbf{protons} & \textbf{neutrons} & \textbf{electrons} & \textbf{charge}\\
\midrule
\ce{^{27}_{13}Al} & \answerblank[1cm]{13} & \answerblank[1cm]{14} &
  \answerblank[1cm]{13} & \answerblank[1cm]{0}\\
\ce{^{19}_{9}F^-} & \answerblank[1cm]{9} & \answerblank[1cm]{10} &
  \answerblank[1cm]{10} & \answerblank[1cm]{$1-$}\\
\ce{^{65}_{30}Zn^{2+}} & \answerblank[1cm]{30} & \answerblank[1cm]{35} &
  \answerblank[1cm]{28} & \answerblank[1cm]{$2+$}\\
? & 20 & 20 & 18 & \answerblank[2.4cm]{\ce{^{40}Ca^{2+}}}\\
\bottomrule
\end{tabular}

\segment{INSTRUCTION B}{20}{Charges from the table}

\notesectionz{Predictable ion charges}{2.6--2.7}
\practice{1}

Main-group metals lose electrons; nonmetals gain them:

\renewcommand{\arraystretch}{1.3}
\noindent\begin{tabular}{@{}lccccc@{}}
\toprule
\textbf{Group} & 1A & 2A & 3A & 6A & 7A\\
\midrule
\textbf{Typical ion} & \answerblank[1.2cm]{$1+$} &
  \answerblank[1.2cm]{$2+$} & \answerblank[1.2cm]{$3+$} &
  \answerblank[1.2cm]{$2-$} & \answerblank[1.2cm]{$1-$}\\
\bottomrule
\end{tabular}

Transition metals are the exception --- most form
\answerblank[4.5cm]{more than one charge}, which is why their names need a
\answerblank[3.4cm]{Roman numeral}.

Memorize the exceptions that never take a numeral: \ce{Ag+}, \ce{Zn^2+},
\ce{Cd^2+}, and \ce{Al^3+}.

\segment{APPLICATION}{20}{Build formulas from charges}

\begin{enumerate}[leftmargin=*,itemsep=0.35em]
  \item Potassium + sulfur: \answerblank[2.5cm]{\ce{K2S}}
  \item Aluminum + oxygen: \answerblank[2.5cm]{\ce{Al2O3}}
  \item Barium + nitrogen: \answerblank[2.5cm]{\ce{Ba3N2}}
  \item An ion has 34 protons and 36 electrons. Identify it fully.
        \answerblank[3.5cm]{\ce{Se^2-}}
\end{enumerate}

\begin{exitticket}
An element in group 2A forms an ion with 10 electrons. Identify the element
and the ion. \answerlines[2]{Loses 2 electrons $\to$ neutral atom had 12
$\to$ magnesium; the ion is \ce{Mg^2+}.}
\end{exitticket}

% =====================================================================
\blocklesson{Nomenclature}{Zumdahl \S2.8}

\begin{objectives}
  \item Name and write formulas for ionic compounds (Type I and II),
        polyatomic-ion compounds, binary covalent compounds, and acids.
\end{objectives}

\reading{Zumdahl \S2.8}{87--98}

\begin{warmup}[3]
  \item Charge on the aluminum ion: \answerblank[2cm]{$3+$}
  \item Neutrons in \ce{^{63}Cu}: \answerblank[2cm]{34}
  \item Formula: magnesium + chlorine: \answerblank[2.5cm]{\ce{MgCl2}}
\end{warmup}

\segment{INSTRUCTION A}{25}{The decision tree}

\notesectionz{Which naming rules apply?}{2.8}
\practice{1}

Ask two questions in order:

\begin{enumerate}[leftmargin=*,itemsep=0.2em]
  \item Does the formula start with \ce{H} (in water)? $\to$ it's an
        \answerblank[1.6cm]{acid}.
  \item Otherwise, is the first element a metal?
        \begin{itemize}[leftmargin=1.4em,itemsep=0.1em]
          \item \textbf{Yes} $\to$ \answerblank[1.8cm]{ionic}. Does the metal
                have a variable charge (transition metals, Pb, Sn)?
                Then include a \answerblank[3.4cm]{Roman numeral}.
          \item \textbf{No} (two nonmetals) $\to$
                \answerblank[3.3cm]{binary covalent} --- use
                \answerblank[1.8cm]{prefixes}.
        \end{itemize}
\end{enumerate}

\notesub{Type I ionic --- fixed-charge metal}

Cation name (element name) $+$ anion root $+$
``\answerblank[1.4cm]{-ide}''.\\
\ce{CaS} $=$ \answerblank[4cm]{calcium sulfide};\quad
\ce{Li3N} $=$ \answerblank[4cm]{lithium nitride}.

\notesub{Type II ionic --- variable-charge metal}

Deduce the cation charge from the anion total, then write the numeral.\\
\ce{Fe2O3}: three \ce{O^2-} $=6-$, so each Fe is
\answerblank[1.4cm]{$3+$} $\to$
\answerblank[4cm]{iron(III) oxide}.\\
\ce{CuCl} $=$ \answerblank[4cm]{copper(I) chloride};\quad
\ce{PbCl2} $=$ \answerblank[4cm]{lead(II) chloride}.

\begin{apctrap}
The Roman numeral is the charge on \emph{one} metal ion, never the number of
atoms. \ce{Fe2O3} is iron(\textsc{iii}) oxide, not iron(\textsc{ii}) --- the
subscript 2 tells you how many Fe, the charge balance tells you the numeral.
\end{apctrap}

\segment{GUIDED PRACTICE}{15}{Ionic drill}

\begin{enumerate}[leftmargin=*,itemsep=0.3em]
  \item \ce{CoBr2}: \answerblank[5cm]{cobalt(II) bromide}
  \item \ce{Al2O3}: \answerblank[5cm]{aluminum oxide}
  \item chromium(III) chloride: \answerblank[3cm]{\ce{CrCl3}}
  \item \ce{MnO2}: \answerblank[5cm]{manganese(IV) oxide}
\end{enumerate}

\segment{INSTRUCTION B}{20}{Polyatomics, covalent, acids}

\notesectionz{Polyatomic ions}{2.8}
\practice{1}

Memorize these --- there is no derivation:

\renewcommand{\arraystretch}{1.25}
\noindent\begin{tabular}{@{}llll@{}}
\toprule
\ce{NH4+} ammonium & \ce{NO3-} nitrate & \ce{SO4^2-} sulfate & \ce{PO4^3-} phosphate\\
\ce{OH-} hydroxide & \ce{NO2-} nitrite & \ce{SO3^2-} sulfite & \ce{CO3^2-} carbonate\\
\ce{CN-} cyanide & \ce{ClO3-} chlorate & \ce{HCO3-} hydrogen carbonate & \ce{C2H3O2-} acetate\\
\bottomrule
\end{tabular}

\notesub{The oxyanion ladder (chlorine as the model)}

\ce{ClO-} \answerblank[2.7cm]{hypochlorite} $\to$
\ce{ClO2-} \answerblank[1.8cm]{chlorite} $\to$
\ce{ClO3-} \answerblank[1.8cm]{chlorate} $\to$
\ce{ClO4-} \answerblank[2.4cm]{perchlorate}.

Fewer O $\to$ \emph{-ite}; more O $\to$ \emph{-ate};
\textit{hypo-} is one below \emph{-ite}, \textit{per-} one above
\emph{-ate}. Same ladder works for Br and I.

\notesub{Binary covalent (Type III) --- prefixes}

mono(1), di(2), tri(3), tetra(4), penta(5), hexa(6).
Drop \emph{mono-} on the \answerblank[2.4cm]{first} element only.\\
\ce{N2O5} $=$ \answerblank[5cm]{dinitrogen pentoxide};\quad
\ce{CO} $=$ \answerblank[4cm]{carbon monoxide}.

\notesub{Acids}

\begin{itemize}[leftmargin=*,itemsep=0.2em]
  \item No oxygen: \emph{hydro-}root\emph{-ic acid}. \ce{HCl(aq)} $=$
        \answerblank[4cm]{hydrochloric acid}.
  \item Anion ends \emph{-ate} $\to$ acid ends
        \answerblank[1.6cm]{-ic}. \ce{H2SO4} $=$
        \answerblank[3.4cm]{sulfuric acid}.
  \item Anion ends \emph{-ite} $\to$ acid ends
        \answerblank[1.8cm]{-ous}. \ce{H2SO3} $=$
        \answerblank[3.4cm]{sulfurous acid}.
\end{itemize}

\segment{APPLICATION}{20}{Mixed naming --- the real test}

Name or write the formula. Identify the type first.
\begin{enumerate}[leftmargin=*,itemsep=0.35em]
  \item \ce{Fe(NO3)3}: \answerblank[5.5cm]{iron(III) nitrate}
  \item \ce{P4O10}: \answerblank[5.5cm]{tetraphosphorus decoxide}
  \item \ce{HNO2(aq)}: \answerblank[5.5cm]{nitrous acid}
  \item sodium hydrogen carbonate: \answerblank[3cm]{\ce{NaHCO3}}
  \item cesium perchlorate: \answerblank[3cm]{\ce{CsClO4}}
  \item ammonium sulfate: \answerblank[3cm]{\ce{(NH4)2SO4}}
\end{enumerate}

\begin{exitticket}
\ce{Cu(NO3)2} and \ce{CuNO3} --- name both, and explain what forces the
difference. \answerlines[2]{Copper(II) nitrate and copper(I) nitrate. Each
\ce{NO3-} is $1-$, so charge balance requires \ce{Cu^2+} in the first and
\ce{Cu+} in the second.}
\end{exitticket}

\end{document}
